\section{序論}
11月22日朝の観測では、数学の難問に対するAI協働の証明主張と形式検証、ローカル実行可能なオープンウェイトモデル、政策上の利用制限を巡る議論が重なった。モデル性能の拡大だけでなく、成果の検証方法や配布・利用の条件が焦点となっている。

音声APIの課金体系、エージェントのランタイム保護、モデル実装基盤の更新も話題となった。以下はGrokが収集した12件を記載順に整理したもので、一次資料の独立再検証は行っていない。観測窓終了後の投稿と未確認の主張は明示して扱う。
